\documentclass[10pt,a4paper]{article}

%-------------------------------------------------------------------------------
% Packages
%-------------------------------------------------------------------------------
\usepackage[
    top=0.7in,
    bottom=0.7in,
    left=0.8in,
    right=0.8in
]{geometry}

\usepackage[T1]{fontenc}
\usepackage{lmodern}
\usepackage{microtype}
\usepackage{xcolor}
\usepackage{listings}
\usepackage{titlesec}
\usepackage{fancyhdr}
\usepackage{enumitem}
\usepackage{parskip}
\usepackage{hyperref}

%-------------------------------------------------------------------------------
% General Formatting
%-------------------------------------------------------------------------------
\setlength{\parindent}{0pt}
\setlength{\parskip}{6pt}

\hypersetup{
    colorlinks=false,
    hidelinks,
    pdftitle={MTRX3760 Lab 1 Report}
}

%-------------------------------------------------------------------------------
% Colours
%-------------------------------------------------------------------------------
\definecolor{lightgray}{RGB}{245,245,245}
\definecolor{darkgray}{RGB}{70,70,70}

%-------------------------------------------------------------------------------
% Section Formatting
%-------------------------------------------------------------------------------
\titleformat{\section}
  {\Large\bfseries}
  {}
  {0pt}
  {}

\titleformat{\subsection}
  {\large\bfseries}
  {}
  {0pt}
  {}

\titleformat{\subsubsection}
  {\normalsize\bfseries}
  {}
  {0pt}
  {}

\titlespacing*{\section}{0pt}{20pt}{8pt}
\titlespacing*{\subsection}{0pt}{14pt}{5pt}
\titlespacing*{\subsubsection}{0pt}{10pt}{4pt}

%-------------------------------------------------------------------------------
% Header / Footer
%-------------------------------------------------------------------------------
\pagestyle{fancy}
\fancyhf{}
\fancyhead[L]{MTRX3760}
\fancyhead[R]{Lab 1}
\fancyfoot[C]{\thepage}

\renewcommand{\headrulewidth}{0.4pt}
\renewcommand{\footrulewidth}{0pt}

%-------------------------------------------------------------------------------
% Program Output Style
%-------------------------------------------------------------------------------
\lstdefinestyle{output}{
    basicstyle=\ttfamily\small,
    backgroundcolor=\color{lightgray},
    frame=single,
    rulecolor=\color{darkgray},
    framerule=0.4pt,
    xleftmargin=8pt,
    xrightmargin=8pt,
    framexleftmargin=6pt,
    framexrightmargin=6pt,
    aboveskip=8pt,
    belowskip=10pt,
    showstringspaces=false,
    columns=fullflexible
}

%-------------------------------------------------------------------------------
% Question Command
%-------------------------------------------------------------------------------
\newcommand{\question}[1]{%
    \vspace{5pt}
    \textbf{#1}\par
    \vspace{2pt}
}

%-------------------------------------------------------------------------------
% Document
%-------------------------------------------------------------------------------
\begin{document}

%-------------------------------------------------------------------------------
% Title Page
%-------------------------------------------------------------------------------
\begin{titlepage}
    \centering

    \vspace*{2.2cm}

    {\Large University of Sydney\par}

    \vspace{1cm}

    {\LARGE\bfseries MTRX3760\par}

    \vspace{0.3cm}

    {\Huge\bfseries Lab 1 Report\par}

    \vspace{1.5cm}

    \rule{0.75\textwidth}{0.5pt}

    \vspace{1.5cm}

    \begin{tabular}{rl}
        \textbf{Student SID 1:} & 530487229 \\[0.4cm]
        \textbf{Student SID 2:} & 540821501 \\[0.4cm]
        \textbf{Practical Section: } & \texttt{[Monday 3pm]} \\[0.4cm]
    \end{tabular}

    \vfill

    {\large Introduction to Object-Oriented Design in C++}

    \vspace{1cm}

\end{titlepage}

\tableofcontents
\newpage

%-------------------------------------------------------------------------------
% A1
%-------------------------------------------------------------------------------
\section{A1 -- Robot}

\subsection{Function-Based}

\subsubsection{Program Output}

\begin{lstlisting}[style=output]
Left motor 0.8, Right motor 0.2
Left motor 0.55, Right motor 0.45
Left motor 0.3, Right motor 0.7
Left motor 0, Right motor 0.5
\end{lstlisting}

\subsection{Object-Based}

\subsubsection{Program Output}

\begin{lstlisting}[style=output]
Left motor 0.8, Right motor 0.2
Left motor 0.55, Right motor 0.45
Left motor 0.3, Right motor 0.7
Left motor 0, Right motor 0.5
\end{lstlisting}

\subsection{Questions}

\question{1. When you added the battery, what changed in \texttt{main()} in each version?}

In the function-based version, \texttt{main()} declares a battery,
initialise it with \texttt{InitBattery()}, and pass a pointer to it into
\texttt{UpdateRobot()} on every cycle, alongside the sensor, controller,
and motors it already passed. In the object-based version, \texttt{main()}
did not change at all, because the battery was added as a new member
inside \texttt{CRobot} and is managed entirely within \texttt{CRobot::Update()}.

\question{2. In the function-based version, what does \texttt{main()} have to know about the robot's parts? What does it know in the object-based version?}

In the function-based version, \texttt{main()} has to know about every
individual part i.e., the sensor, controller, both motors, and now the
battery, and passes pointers to each of them into the free functions that
operate on them. In the object-based version, \texttt{main()} only needs
to know about the single \texttt{CRobot} object and calls its
\texttt{Update()} and \texttt{Report()} functions; it has no knowledge of
the sensor, controller, motors, or battery that \texttt{CRobot} contains.

\question{3. What is the name of the object-oriented principle that produced this saving?}

Encapsulation. The robot's internal parts and the logic that coordinates
them are kept inside \texttt{CRobot}, so adding a new part (the battery)
only required changing \texttt{CRobot} itself, not any code that uses it.

%-------------------------------------------------------------------------------
% A2
%-------------------------------------------------------------------------------
\newpage
\section{A2 -- Oven}

\subsection{Private Version}

\subsubsection{Program Output}

\begin{lstlisting}[style=output]
Reflow oven is at 25C
Curing oven is at 28C
\end{lstlisting}

\subsection{Public Version}

\subsubsection{Program Output}

\begin{lstlisting}[style=output]
Reflow oven is at 25C
Curing oven is at 28C
\end{lstlisting}

\subsection{Questions}

\question{1. Which version required more changes? Name the functions or lines you had to edit in each.}

The public-data version required more changes, and they were spread across
\texttt{main()}: the starting temperatures changed from \texttt{20} to
\texttt{200}, the warming step changed from \texttt{+= 1} to \texttt{+= 10},
the overheating limit changed from \texttt{250} to \texttt{2500}, and the
printed value had to be divided by \texttt{10} at the point it was printed.
The private version required changes only inside \texttt{COven}: the
constructor's starting value (\texttt{20} to \texttt{200}), \texttt{WarmUp()}
(\texttt{+= 1} to \texttt{+= 10}), \texttt{IsOverheating()}'s limit
(\texttt{250} to \texttt{2500}), and \texttt{Report()} (dividing by \texttt{10}
before printing). \texttt{main()} itself needed no changes at all.

\question{2. What is the name of the object-oriented principle that kept the change contained in one version?}

Encapsulation. Because the temperature is private in \texttt{COven}, the
details of how it is stored and scaled can change inside the class without
any other code needing to know or change.

\question{3. Did \texttt{main()} need to change in the public version? In the private version? Explain the difference.}

Yes, \texttt{main()} had to change in the public version, because it reads
and writes \texttt{mTemperatureC} directly and therefore has to know the
new scale. In the private version, \texttt{main()} did not need to change,
because it only calls \texttt{WarmUp()}, \texttt{IsOverheating()}, and
\texttt{Report()} functions whose interface stayed the same even though
their internal implementation changed.

%-------------------------------------------------------------------------------
% A3
%-------------------------------------------------------------------------------
\newpage
\section{A3 -- Clock}

\subsection{Standalone Version}

\subsubsection{Program Output}

\begin{lstlisting}[style=output]
Kitchen 07:00
Alarm 07:02
\end{lstlisting}

\subsection{Inheritance Version}

\subsubsection{Program Output}

\begin{lstlisting}[style=output]
Kitchen 07:00
Alarm 07:02
\end{lstlisting}

\subsection{Questions}

\question{1. In the standalone version, which parts of the clock's behaviour did you have to repeat?}

The standalone \texttt{CAlarmClock} had to repeat the plain clock's time
data and behaviour entirely: it stores its own name, start time, and
current time, and duplicates \texttt{Tick()} to advance the time and
\texttt{Report()} to format and print it, even though \texttt{CClock}
already implements both.

\question{2. In the inheritance version, what did you add, and what did you get without writing it?}

\texttt{CAlarmClock} added only what an alarm needs on top of a clock: an
\texttt{mAlarmMinutes} member, \texttt{SetAlarm()}, and \texttt{IsRinging()}.
Everything else, like \texttt{Tick()}, \texttt{GetTime()}, \texttt{Report()},
\texttt{Reset()}, and the constructor's minute-tracking was inherited
from \texttt{CClock} and did not need to be written or duplicated.

\question{3. What is the name of the object-oriented principle that let the second version reuse the clock's behaviour?}

Inheritance. \texttt{CAlarmClock} reuses \texttt{CClock}'s existing
behaviour by deriving from it, adding only its own alarm-specific state
and functions.

%-------------------------------------------------------------------------------
% A4
%-------------------------------------------------------------------------------
\newpage
\section{A4 -- Power Budget}

\subsection{Power Budget Version 1}

\subsubsection{Program Output}

\begin{lstlisting}[style=output]
DriveMotor: 60 W
StatusLed: 0 W
Heater: 60 W
Total: 120 W
\end{lstlisting}

\subsection{Power Budget Version 2}

\subsubsection{Program Output}

\begin{lstlisting}[style=output]
DriveMotor: 60 W
StatusLed: 0 W
Heater: 60 W
Total: 120 W
\end{lstlisting}

\subsection{Questions}

\question{1. What did you have to change in each version to add the heater? Answer for both versions separately.}

In \texttt{power\_budget1.cpp}, a \texttt{CHeater} class was added as a
subclass of \texttt{CDevice}, without modifying \texttt{CDevice}, with a
heat setting, \texttt{SetHeat()}, and its own \texttt{PowerDraw()}. In
\texttt{main()}, a heater object was created, turned on, and given a heat
setting, then a new three-line block was written by hand to get its
\texttt{PowerDraw()}, print it, and add it to the running total,
copying the pattern already used for the motor and the LED.

In \texttt{power\_budget2.cpp}, the same \texttt{CHeater} class was added,
also without modifying \texttt{CDevice}. The only change needed in
\texttt{main()} was constructing the heater and adding a pointer to it in
the \texttt{devices} array. No new reporting or totalling code was needed,
because the existing loop already handles every \texttt{CDevice*} in the
array polymorphically.

\question{2. Which version was easier to extend, and what specifically made it easier?}

\texttt{power\_budget2.cpp} was easier to extend. Because every device is
stored and processed as a \texttt{CDevice*} through one loop, adding the
heater only meant creating it and adding it to the array, the loop
picked it up automatically. \texttt{power\_budget1.cpp} required writing
an entire new block in \texttt{main()} to get, print, and total the
heater's draw, duplicating the pattern already present for the other two
devices.

\question{3. What is the name of the object-oriented principle behind that difference?}

Polymorphism. Because \texttt{CMotor}, \texttt{CLed}, and \texttt{CHeater}
all implement the same \texttt{PowerDraw()} interface declared in
\texttt{CDevice}, the loop in \texttt{main()} can call it on any of them
through a \texttt{CDevice*} without knowing which concrete class it is
pointing to.





%-------------------------------------------------------------------------------
% A5
%-------------------------------------------------------------------------------
\newpage
\section{A5 -- Robot Controller}

\subsection{Program Output}

\begin{lstlisting}[style=output]
Cycle 1
Drive motor: current 0.2, target 0.8
Front line offset: 2

Cycle 2
Drive motor: current 0.4, target 0.8
Front line offset: 1

Cycle 3
Drive motor: current 0.6, target 0.8
Front line offset: -1

Cycle 4
Drive motor: current 0.8, target 0.8
Front line offset: -2
\end{lstlisting}

\subsection{Questions}
\question{1. What did Partner B's controller need to know about the subsystems? What did Partner A's subsystems need to know about the controller?}

Partner B's controller only needed each subsystem to provide the common
\texttt{CSubsystem} interface: \texttt{Update()} and \texttt{Report()}.
Partner A's subsystems did not need to know about \texttt{CRobotController},
they only inherited from \texttt{CSubsystem} and implemented that interface.

\question{2. What did you have to change, adapt, or correct to make the two halves work together?}

We made both subsystem classes inherit from \texttt{CSubsystem} and ensured
their \texttt{Update()} and \texttt{Report() const} function signatures
matched the base-class interface. In \texttt{main.cpp}, we created the
subsystems, configured the motor target speed, and added pointers to both
subsystems to the controller.

\question{3. Did you design the testing before or after the subsystems and integrated system? In future, which option do you think would be the most time-efficient?}

We designed the integration test after implementing the subsystems and
controller, using \texttt{main.cpp} to run four cycles and print their state.
In future, defining expected test output before implementation would be more
time-efficient because it would make integration errors easier to identify.

\question{4. In which places does your design, including test code, use the following concepts? Be specific, point to specific parts of the program:}

\textbf{a. Modularity:} Each class is in its own header and implementation
file, including \texttt{CDriveMotor}, \texttt{CLineDetector},
\texttt{CSubsystem}, and \texttt{CRobotController}.

\textbf{b. Encapsulation:} The motor's speeds, detector's line offset, and
controller's subsystem array are private member data, accessed only through
their public member functions.

\textbf{c. Inheritance:} \texttt{CDriveMotor} and \texttt{CLineDetector}
both publicly inherit from the abstract \texttt{CSubsystem} base class.

\textbf{d. Polymorphism:} \texttt{CRobotController} stores
\texttt{CSubsystem*} pointers and calls the virtual \texttt{Update()} and
\texttt{Report()} functions in its loops, so it can operate on either type
without knowing its concrete class.

%-------------------------------------------------------------------------------
% Self-Assessment
%-------------------------------------------------------------------------------
\newpage
\section{Self-Assessment}

\texttt{Without bonus}

\begin{center}
\begin{tabular}{ll}
A1 & \_\_\_ / 5 \\
A2 & \_\_\_ / 5 \\
A3 & \_\_\_ / 5 \\
A4 & \_\_\_ / 5 \\
A5 & \_\_\_ / 30 \\
Design & \_\_\_ / 25 \\
Code & \_\_\_ / 25 \\
\textbf{Total} & \textbf{\_\_\_ / 100} \\
\end{tabular}
\end{center}

%-------------------------------------------------------------------------------
% Code Appendix
%-------------------------------------------------------------------------------
\clearpage
\section{Code Appendix}

\lstinputlisting[style=output,language=C++,breaklines=true,
caption={A1: Function-Based Robot}]
{A1.Robot/robot_function_based.cpp}

\newpage
\lstinputlisting[style=output,language=C++,breaklines=true,
caption={A1: Object-Based Robot}]
{A1.Robot/robot_object_based.cpp}
\newpage
\lstinputlisting[style=output,language=C++,breaklines=true,
caption={A2: Public Oven}]
{A2.Oven/oven_public.cpp}
\newpage
\lstinputlisting[style=output,language=C++,breaklines=true,
caption={A2: Private Oven}]
{A2.Oven/oven_private.cpp}
\newpage
\lstinputlisting[style=output,language=C++,breaklines=true,
caption={A3: Standalone Clock}]
{A3.Clock/clock_standalone.cpp}
\newpage
\lstinputlisting[style=output,language=C++,breaklines=true,
caption={A3: Inheritance Clock}]
{A3.Clock/clock_inherit.cpp}
\newpage
\lstinputlisting[style=output,language=C++,breaklines=true,
caption={A4: Power Budget Version 1}]
{A4.PowerBudget/power_budget1.cpp}
\newpage
\lstinputlisting[style=output,language=C++,breaklines=true,
caption={A4: Power Budget Version 2}]
{A4.PowerBudget/power_budget2.cpp}
\newpage
\lstinputlisting[style=output,language=C++,breaklines=true,
caption={A5: CSubsystem Header}]
{A5.RobotController/CSubsystem.h}
\newpage
\lstinputlisting[style=output,language=C++,breaklines=true,
caption={A5: CSubsystem Implementation}]
{A5.RobotController/CSubsystem.cpp}
\newpage
\lstinputlisting[style=output,language=C++,breaklines=true,
caption={A5: CDriveMotor Header}]
{A5.RobotController/CDriveMotor.h}
\newpage
\lstinputlisting[style=output,language=C++,breaklines=true,
caption={A5: CDriveMotor Implementation}]
{A5.RobotController/CDriveMotor.cpp}
\newpage
\lstinputlisting[style=output,language=C++,breaklines=true,
caption={A5: CLineDetector Header}]
{A5.RobotController/CLineDetector.h}
\newpage
\lstinputlisting[style=output,language=C++,breaklines=true,
caption={A5: CLineDetector Implementation}]
{A5.RobotController/CLineDetector.cpp}
\newpage
\lstinputlisting[style=output,language=C++,breaklines=true,
caption={A5: CRobotController Header}]
{A5.RobotController/CRobotController.h}
\newpage
\lstinputlisting[style=output,language=C++,breaklines=true,
caption={A5: CRobotController Implementation}]
{A5.RobotController/CRobotController.cpp}
\newpage
\lstinputlisting[style=output,language=C++,breaklines=true,
caption={A5: Integration Program}]
{A5.RobotController/main.cpp}

\end{document}
